% !TEX program = pdflatex
% !TEX option  = -synctex=1 -interaction=nonstopmode -output-directory=build

\newif\ifhandout
\handouttrue            % <- cambia a \handoutfalse para presentación normal
%\handoutfalse

\ifhandout
\documentclass[10pt,handout,aspectratio=169]{beamer}
\else
\documentclass[10pt,aspectratio=169]{beamer}
\fi

% ------------------------------------------------
% PAQUETES
% ------------------------------------------------
\usepackage[spanish,es-noshorthands]{babel}
\usepackage[T1]{fontenc}
\usepackage[utf8]{inputenc}
\usepackage{mathpazo}
\usepackage{amsmath,amssymb,amsfonts}
\usepackage{hyperref}
\usepackage{tikz}
\usepackage{graphicx}
\usepackage{booktabs}
\usepackage{array}
\usepackage{multicol}
\usepackage{ragged2e}

% ------------------------------------------------
% TEMA Y ESTILO
% ------------------------------------------------
\usetheme{Madrid}
\usecolortheme[named=teal]{structure}
\setbeamertemplate{navigation symbols}{}
\setbeamertemplate{itemize items}[circle]
\setbeamertemplate{enumerate items}[default]

\definecolor{verdeagro}{RGB}{0,153,117}
\definecolor{gristexto}{RGB}{50,50,50}

\setbeamercolor{title}{fg=verdeagro}
\setbeamercolor{frametitle}{fg=white,bg=verdeagro}
\setbeamercolor{structure}{fg=verdeagro}
\setbeamercolor{block title}{bg=verdeagro!90,fg=white}
\setbeamercolor{block body}{bg=verdeagro!6,fg=black}
\setbeamercolor{alerted text}{fg=red!75!black}
\setbeamercolor{normal text}{fg=gristexto,bg=white}

\setbeamertemplate{footline}{
	\leavevmode%
	\hbox{%
		\begin{beamercolorbox}[wd=.75\paperwidth,ht=2.5ex,dp=1ex,leftskip=1em]{author in head/foot}%
			\usebeamerfont{author in head/foot}Matemáticas
		\end{beamercolorbox}%
		\begin{beamercolorbox}[wd=.25\paperwidth,ht=2.5ex,dp=1ex,rightskip=1em plus1fil]{date in head/foot}%
			\usebeamerfont{date in head/foot}\insertframenumber/\inserttotalframenumber
	\end{beamercolorbox}}%
	\vskip0pt%
}

% ------------------------------------------------
% ENTORNOS CON OVERLAY
% ------------------------------------------------
\newenvironment{stepenumerate}{\begin{enumerate}[<+->]}{\end{enumerate}}
\newenvironment{stepitemize}{\begin{itemize}[<+->]}{\end{itemize}}
\newenvironment{stepenumeratewithalert}{\begin{enumerate}[<+-| alert@+>]}{\end{enumerate}}
\newenvironment{stepitemizewithalert}{\begin{itemize}[<+-| alert@+>]}{\end{itemize}}

% ------------------------------------------------
% DATOS
% ------------------------------------------------
\title[Matemáticas]{Matemáticas}
\subtitle{Clase 2 \\ Expresiones algebraicas, exponentes y radicales}
\author[Luis González Alcaino]{Luis González Alcaino\\Magíster en Matemática}
\institute[UCM]{Universidad Católica del Maule\\Curicó}
\date{Versión actualizada}

% ------------------------------------------------
% DOCUMENTO
% ------------------------------------------------
\begin{document}
	
	% ------------------------------------------------
	\begin{frame}
		\titlepage
	\end{frame}
	
	% ------------------------------------------------
	\section{Matemática}
	
	\begin{frame}{Contenidos de la clase}
		\begin{multicols}{2}
			\begin{itemize}
				\item Álgebra
				\item Término algebraico
				\item Expresiones algebraicas
				\item Clasificación de expresiones algebraicas
				\item Polinomios
				\item Valor numérico de una expresión algebraica
				\item Términos semejantes
				\item Reducción de términos semejantes
				\item Eliminación de paréntesis
				\item Notación exponencial
				\item Leyes de los exponentes
				\item Exponentes racionales
				\item Ejercicios
				\item Lecturas
			\end{itemize}
		\end{multicols}
	\end{frame}
	
	% ------------------------------------------------
	\subsection{Álgebra}
	
	\begin{frame}{Álgebra}
		\justifying
		Es la rama de la matemática que estudia la cantidad considerada del modo más general posible. Puede definirse como la generalización y extensión de la aritmética.
		
		\bigskip
		
		A diferencia de la aritmética elemental, que trata de los números y las operaciones fundamentales, en álgebra se introducen letras para representar variables o cantidades desconocidas. Las expresiones así formadas se llaman \emph{fórmulas algebraicas} y expresan reglas o principios generales.
		
		\bigskip
		
		En síntesis, el álgebra constituye una generalización de la aritmética, ya que mediante la combinación de números y letras podemos dar solución a situaciones más complejas.
	\end{frame}
	
	% ------------------------------------------------
	\subsection{Término algebraico}
	
	\begin{frame}{Término algebraico}
		Un término algebraico es una expresión elemental donde aparecen solamente operaciones de multiplicación y división entre números y letras. 
		
		\bigskip
		
		El número se llama \textbf{coeficiente numérico} y las letras conforman la \textbf{parte literal}. Tanto el número como cada letra pueden estar elevados a una potencia.
		
		\begin{block}{Ejemplo}
			Término algebraico: \[
			-9a^{3}x^{5}y^{2}
			\]
			
			\begin{itemize}
				\item $-9$ es el coeficiente numérico.
				\item $a^{3}x^{5}y^{2}$ es la parte literal.
				\item $3$, $5$ y $2$ son exponentes.
			\end{itemize}
		\end{block}
	\end{frame}
	
	% ------------------------------------------------
	\subsection{Expresiones algebraicas}
	
	\begin{frame}{Expresión algebraica}
		Cuando se combinan términos algebraicos mediante operaciones de suma, resta, multiplicación, división, potenciación o extracción de raíces, la expresión resultante se llama \textbf{expresión algebraica}.
		
		\begin{block}{Ejemplos de expresiones algebraicas}
			\begin{enumerate}
				\item $5ax^{3}-bx^{2}+5$
				\item $15-4\sqrt{x}+\dfrac{4y}{5x}$
				\item $\dfrac{(x+2z)^{2}-5xz}{x+z^{2}}-10xz^{3}$
			\end{enumerate}
		\end{block}
	\end{frame}
	
	% ------------------------------------------------
	\subsection{Clasificación de expresiones algebraicas}
	
	\begin{frame}{Clasificación de expresiones algebraicas}
		Las expresiones algebraicas que tienen exactamente un término se denominan \textbf{monomios}. Aquellas que tienen exactamente dos términos son \textbf{binomios}, y las que tienen exactamente tres términos son \textbf{trinomios}. Las que tienen más de un término se llaman en general \textbf{multinomios}.
		
		\begin{block}{Ejemplo}
			\begin{enumerate}
				\item $4a^{3}b^{2}$ es un monomio.
				\item $5x^{3}-3x$ es un binomio.
				\item $7x^{2}+\sqrt{y^{5}}-15$ es un trinomio.
				\item $2\sqrt{x}-3x^{2}-\dfrac{3}{x}+4x-5$ es un multinomio.
			\end{enumerate}
		\end{block}
	\end{frame}
	
	% ------------------------------------------------
	\subsection{Polinomios}
	
	\begin{frame}{Polinomios}
		\justifying
		Un \textbf{polinomio} es una expresión algebraica formada por un número finito de variables y constantes, usando únicamente las operaciones de suma, resta y multiplicación, junto con exponentes enteros no negativos.
		
		\bigskip
		
		Más precisamente, es una combinación lineal de productos de potencias enteras positivas de una o varias variables.
		
		\begin{block}{Ejemplos}
			\begin{enumerate}
				\item $3x^{2}+3x-1$ es un polinomio en una variable, de grado $2$.
				\item $5xy^{3}-2xy+x-y+1$ es un polinomio en dos variables, de grado $4$.
				\item $x^{-2}+3x-4$ no es un polinomio.
				\item $\dfrac{x^{3}-7x^{2}+1}{x+y^{2}}-2xy^{3}$ no es un polinomio.
			\end{enumerate}
		\end{block}
	\end{frame}
	
	% ------------------------------------------------
	\subsection{Valor numérico de una expresión algebraica}
	
	\begin{frame}{Valor numérico de una expresión algebraica}
		El valor numérico de una expresión algebraica es el valor que toma la expresión cuando se asignan valores numéricos a sus variables.
		
		\begin{block}{Ejemplos}
			\begin{enumerate}
				\item Para $3x+2x^{2}-5$, cuando $x=\dfrac{1}{2}$:
				\[
				3\left(\frac12\right)+2\left(\frac12\right)^2-5
				=\frac32+\frac12-5=-3
				\]
				
				\item Para $9x^{3}y+4x^{2}y^{3}-3xy$, cuando $x=\dfrac13$ e $y=-\dfrac12$:
				\[
				9\left(\frac13\right)^3\left(-\frac12\right)
				+4\left(\frac13\right)^2\left(-\frac12\right)^3
				-3\left(\frac13\right)\left(-\frac12\right)
				=\frac{5}{18}
				\]
				
				\item Hallar el área de un trapecio usando
				\[
				A=\frac{1}{2}h(b+c)
				\]
				si $h=6$, $b=12$ y $c=8$:
				\[
				A=\frac{1}{2}(6)(12+8)=3(20)=60\text{ m}^2
				\]
			\end{enumerate}
		\end{block}
	\end{frame}
	
	% ------------------------------------------------
	\subsection{Términos semejantes}
	
	\begin{frame}{Términos semejantes}
		Los términos semejantes son aquellos que tienen exactamente la misma parte literal, es decir, las mismas letras elevadas a los mismos exponentes, y difieren solo en el coeficiente numérico.
		
		\begin{block}{Ejemplos}
			\begin{enumerate}
				\item Son términos semejantes: $3x^{2}y$, $-7x^{2}y$, $9x^{2}y$.
				\item No son términos semejantes: $2x^{2}y^{3}$, $-4x^{2}y$, $12x^{2}y^{4}$.
			\end{enumerate}
		\end{block}
	\end{frame}
	
	% ------------------------------------------------
	\subsection{Reducción de términos semejantes}
	
	\begin{frame}{Reducción de términos semejantes}
		Es la operación cuyo objetivo es convertir en un solo término dos o más términos semejantes.
		
		\begin{itemize}
			\item \textbf{Mismo signo}
			\[
			3x^{2}y+5x^{2}y=8x^{2}y
			\]
			\[
			-7a^{3}b-12a^{3}b=-19a^{3}b
			\]
			
			\item \textbf{Distinto signo}
			\[
			3xyz-4xyz=-xyz
			\]
			\[
			-6ab^{4}+9ab^{4}=3ab^{4}
			\]
		\end{itemize}
	\end{frame}
	
	% ------------------------------------------------
	\subsection{Eliminación de paréntesis}
	
	\begin{frame}{Eliminación de paréntesis}
		El uso de paréntesis en álgebra es muy frecuente. Para simplificar expresiones algebraicas, es necesario eliminarlos correctamente.
		
		\begin{itemize}
			\item Si delante del paréntesis hay un signo $+$, se eliminan sin cambiar signos.
			\item Si delante del paréntesis hay un signo $-$, se eliminan cambiando el signo de todos los términos interiores.
			\item Si hay varios paréntesis, se trabaja desde el interior hacia el exterior.
		\end{itemize}
		
		\begin{block}{Ejemplos}
			\begin{enumerate}
				\item 
				\[
				7x+(-5y+6z)-(8z-3y+4x)
				\]
				\[
				=7x-5y+6z-8z+3y-4x=3x-2y-2z
				\]
				
				\item
				\[
				3ab-\{3a-(-5ab+8a)-2a\}
				\]
				\[
				=3ab-\{3a+5ab-8a-2a\}
				=3ab-\{-7a+5ab\}
				\]
				\[
				=3ab+7a-5ab=7a-2ab
				\]
			\end{enumerate}
		\end{block}
	\end{frame}
	
	% ------------------------------------------------
	\subsection{Exponentes enteros}
	
	\begin{frame}{Notación exponencial}
		Si $a$ es un número real y $n$ es un entero positivo, entonces la potencia $n$-ésima de $a$ es
		\[
		a^{n}=\underbrace{a\cdot a\cdot \ldots \cdot a}_{n\text{ factores}}
		\]
		
		El número $a$ se denomina \textbf{base} y $n$ el \textbf{exponente}.
		
		\medskip
		
		Si $a\neq 0$ y $n$ es entero positivo:
		\[
		a^{0}=1
		\qquad\text{y}\qquad
		a^{-n}=\frac{1}{a^{n}}
		\]
		
		\begin{block}{Ejemplos}
			\begin{itemize}
				\item $3^{4}=81$
				\item $(-2)^{6}=64$
				\item $-5^{3}=-(5^{3})=-125$
				\item $-4^{-2}=-\dfrac{1}{4^{2}}=-\dfrac{1}{16}$
				\item $7^{0}=1$
			\end{itemize}
		\end{block}
	\end{frame}
	
	% ------------------------------------------------
	\subsection{Leyes de los exponentes}
	
	\begin{frame}{Leyes de los exponentes}
		\small
		\begin{tabular}{>{\raggedright\arraybackslash}p{0.32\textwidth} >{\raggedright\arraybackslash}p{0.24\textwidth} >{\raggedright\arraybackslash}p{0.36\textwidth}}
			\toprule
			\textbf{Ley} & \textbf{Ejemplo} & \textbf{Descripción} \\
			\midrule
			$a^{m}\cdot a^{n}=a^{m+n}$ & $3^{2}\cdot 3^{5}=3^{7}$ & Para multiplicar potencias de la misma base, se suman los exponentes. \\[0.4em]
			$\dfrac{a^{m}}{a^{n}}=a^{m-n}$ & $\dfrac{4^{5}}{4^{2}}=4^{3}$ & Para dividir potencias de la misma base, se restan los exponentes. \\[0.4em]
			$(a^{m})^{n}=a^{mn}$ & $(5^{2})^{3}=5^{6}$ & Para elevar una potencia a otra, se multiplican los exponentes. \\[0.4em]
			$(ab)^{n}=a^{n}b^{n}$ & $(3\cdot 4)^{2}=3^{2}\cdot 4^{2}$ & La potencia de un producto se distribuye a cada factor. \\
			\bottomrule
		\end{tabular}
	\end{frame}
	
	\begin{frame}{Leyes de los exponentes}
		\small
		\begin{tabular}{>{\raggedright\arraybackslash}p{0.34\textwidth} >{\raggedright\arraybackslash}p{0.24\textwidth} >{\raggedright\arraybackslash}p{0.34\textwidth}}
			\toprule
			\textbf{Ley} & \textbf{Ejemplo} & \textbf{Descripción} \\
			\midrule
			$\left(\dfrac{a}{b}\right)^n=\dfrac{a^n}{b^n}$ &
			$\left(\dfrac{3}{2}\right)^4=\dfrac{3^4}{2^4}$ &
			La potencia de un cociente se distribuye al numerador y denominador. \\[0.5em]
			
			$\left(\dfrac{a}{b}\right)^{-n}=\left(\dfrac{b}{a}\right)^n$ &
			$\left(\dfrac{3}{2}\right)^{-2}=\left(\dfrac{2}{3}\right)^2$ &
			Una potencia negativa invierte la fracción y cambia el signo del exponente. \\[0.5em]
			
			$\dfrac{a^{-n}}{b^{-m}}=\dfrac{b^m}{a^n}$ &
			$\dfrac{3^{-2}}{2^{-5}}=\dfrac{2^5}{3^2}$ &
			Al mover un factor entre numerador y denominador, cambia el signo del exponente. \\
			\bottomrule
		\end{tabular}
	\end{frame}
	
	% ------------------------------------------------
	\subsection{Exponentes racionales}
	
	\begin{frame}{Exponentes racionales}
		\begin{block}{Raíz $n$-ésima}
			Si $n$ es un entero positivo, la raíz $n$-ésima principal de $a$ se define por:
			\[
			\sqrt[n]{a}=b
			\qquad \Longleftrightarrow \qquad
			b^{n}=a
			\]
			Si $n$ es par, se requiere $a\geq 0$ y $b\geq 0$.
		\end{block}
		
		\begin{block}{Propiedades de las raíces $n$-ésimas}
			\begin{enumerate}
				\item $\sqrt[n]{ab}=\sqrt[n]{a}\sqrt[n]{b}$
				\item $\sqrt[n]{\dfrac{a}{b}}=\dfrac{\sqrt[n]{a}}{\sqrt[n]{b}}$
				\item $\sqrt[m]{\sqrt[n]{a}}=\sqrt[mn]{a}$
				\item $\sqrt[n]{a^{n}}=a$, si $n$ es impar
				\item $\sqrt[n]{a^{n}}=|a|$, si $n$ es par
			\end{enumerate}
		\end{block}
	\end{frame}
	
	\begin{frame}{Exponentes racionales}
		Para cualquier exponente racional $\dfrac{m}{n}$, con $m,n\in\mathbb{Z}$ y $n>0$, se define:
		\[
		a^{\frac{m}{n}}=\sqrt[n]{a^{m}}
		\qquad\text{o equivalentemente}\qquad
		a^{\frac{m}{n}}=\left(\sqrt[n]{a}\right)^{m}
		\]
		
		En particular,
		\[
		a^{\frac{1}{n}}=\sqrt[n]{a}
		\]
		
		\medskip
		
		Si $n$ es par, entonces es necesario que $a\geq 0$.
		
		\bigskip
		
		Con esta definición, las leyes de los exponentes siguen siendo válidas para exponentes racionales.
	\end{frame}
	
	% ------------------------------------------------
	\subsection{Ejercicios}
	
	\begin{frame}{Ejercicios}
		\begin{stepenumerate}
			\item Reduzca los siguientes términos semejantes:
			
			\begin{stepitemize}
				\item $-(12x^{3}y+7xy^{3})-(-4x^{3}y-8xy^{3})+3xy^{3}-(5x^{3}y-8x^{3}y)$
				
				\item $\dfrac{1}{4}ab-\dfrac{3}{2}a^{2}b+\dfrac{2}{5}ab^{2}-\dfrac{1}{3}ab+\dfrac{2}{3}ab^{2}-a^{2}b$
				
				\item $3x-\left(-2x+3y-(4x+y)\right)$
				
				\item $3a-[-(-b+3a)-a+2b-(10b-3a)-3b-2(-b+a)]$
			\end{stepitemize}
			
			\item Si $2a=-5$, $b=-1$, $c=-2$ y $2d=4$, entonces:
			\[
			4a+2c-2b-4d=
			\]
			
			\item Al restar $(2a^{2}-3ab+2b^{2})$ de $(5a^{2}-3ab-b^{2})$, ¿qué resulta?
			
			\item En geometría, la fórmula de Herón es:
			\[
			A=\sqrt{s(s-a)(s-b)(s-c)}, \qquad s=\frac{a+b+c}{2}
			\]
			Calcule el área de un triángulo cuyos lados miden $4$, $10$ y $12$ cm.
		\end{stepenumerate}
	\end{frame}
	
	\begin{frame}{Ejercicios}
		\begin{stepenumerate}
			\item Simplifique las siguientes expresiones algebraicas:
			
			\begin{stepenumerate}
				\item $\left(\dfrac{2x^{2}}{4y^{3}}\right)^{3}$
				
				\item $\dfrac{\left(-a^{3}bc\right)\left(2ab^{-1}c^{2}\right)}{(ab^{2})\left(-a^{-2}b^{3}c^{3}\right)(-a)}$
				
				\item $\dfrac{\left(r^{-2}s^{3}pq\right)^{-3}}{\left(r^{3}s^{2}\right)^{-2}\left(pq^{3}\right)^{-2}}$
				
				\item $\dfrac{\sqrt{\sqrt[3]{x}}}{\sqrt{y}}\div\dfrac{\sqrt[3]{\sqrt{x}}}{\sqrt{xy}}$
				
				\item $\sqrt[3]{32x\sqrt{x^{3}y^{9}}\sqrt{xy^{3}}}$
			\end{stepenumerate}
			
			\item Determine el valor de $n$ en los siguientes enunciados:
			
			\begin{stepenumerate}
				\item $\dfrac{1}{16}=2^{n}$
				
				\item $\dfrac{10^{6}\times 10^{-5}}{10^{-3}}=10^{n}$
				
				\item $\dfrac{1}{\left(\dfrac{\sqrt{2}x^{-2}}{\sqrt{12}x^{3}}\right)^{2}}=6x^{n}$
			\end{stepenumerate}
		\end{stepenumerate}
	\end{frame}
	
	% ------------------------------------------------
	\subsection{Soluciones de los ejercicios}
	
	\begin{frame}{Soluciones de los ejercicios}
		\begin{enumerate}
			\item Reducción de términos semejantes:
			
			\begin{itemize}
				\item $4xy^{3}-5x^{3}y$
				\item $-\dfrac{5}{2}a^{2}b+\dfrac{16}{15}ab^{2}-\dfrac{1}{12}ab$
				\item $9x-2y$
				\item $6a+8b$
			\end{itemize}
			
			\item $-20$
			
			\item 
			\[
			(5a^{2}-3ab-b^{2})-(2a^{2}-3ab+2b^{2})=3a^{2}-3b^{2}
			\]
			
			\item 
			\[
			A=\sqrt{13(13-4)(13-10)(13-12)}=\sqrt{351}\approx 18{,}7\text{ cm}^{2}
			\]
		\end{enumerate}
	\end{frame}
	
	\begin{frame}{Soluciones de los ejercicios}
		\begin{stepenumerate}
			\item Simplificaciones:
			
			\begin{stepenumerate}
				\item $\dfrac{x^{6}}{8y^{9}}$
				\item $-\dfrac{2a^{4}}{b^{5}}$
				\item $\dfrac{q^{3}r^{12}}{ps^{5}}$
				\item $\sqrt{x}$
			\end{stepenumerate}
			
			\item Valor de $n$:
			
			\begin{stepenumerate}
				\item $n=-4$
				\item $n=4$
				\item $n=10$
			\end{stepenumerate}
		\end{stepenumerate}
	\end{frame}
	
	% ------------------------------------------------
	\subsection{Lecturas}
	
	\begin{frame}{Lecturas}
		Páginas 9--14 y 16--26 del libro:
		
		\bigskip
		
		\textbf{Haeussler, Ernest F. \& Paul Jr., Richard S.} (2008).\\
		\emph{Matemáticas para la Administración y Economía}.\\
		Pearson, México.
	\end{frame}
	
	% ------------------------------------------------
	\begin{frame}
		\centering
		\Huge Gracias
	\end{frame}
	
\end{document}